\begin{proposition}[碰撞条件下汉明距离的闭式母函数与子集能量展开]\label{prop:fold-collision-conditional-hamming-mgf}
取 \(\omega,\eta\overset{\mathrm{i.i.d.}}{\sim}\mathrm{Unif}(\Omega_m)\) 并令 \(M:=F_{m+2}\)。
设碰撞事件
\[
\mathcal C_m:=\{\Fold_m(\omega)=\Fold_m(\eta)\}.
\]
记汉明距离
\[
D_m:=\card{\{1\le k\le m:\ \omega_k\neq \eta_k\}}.
\]
对 \(0\le t\le M-1\) 与 \(1\le k\le m\) 定义角变量
\[
\theta_{k,t}:=\frac{\pi t F_{k+1}}{M}.
\]
则对任意复数 \(u\) 成立严格恒等式
\[
\boxed{\;
\EE\bigl[u^{D_m}\mathbf 1_{\mathcal C_m}\bigr]
=\frac1M\sum_{t=0}^{M-1}\prod_{k=1}^m \frac{1+u\cos(2\theta_{k,t})}{2}.\;}
\]
从而碰撞概率满足
\[
\boxed{\;
\mathsf{Col}_m:=\PP(\mathcal C_m)
=\frac1M\sum_{t=0}^{M-1}\prod_{k=1}^m \cos^2(\theta_{k,t}),\;}
\]
并且条件母函数具有闭式
\[
\boxed{\;
\EE\bigl[u^{D_m}\mid \mathcal C_m\bigr]
=\frac{\displaystyle \frac1M\sum_{t=0}^{M-1}\prod_{k=1}^m \frac{1+u\cos(2\theta_{k,t})}{2}}
{\displaystyle \mathsf{Col}_m}.\;}
\]
\par
进一步，对任意子集 \(I\subseteq[m]\) 定义子集能量
\[
E_{m,I}:=\frac1M\sum_{t=0}^{M-1}\Bigl(\prod_{k\in I}\sin^2(\theta_{k,t})\Bigr)\Bigl(\prod_{k\notin I}\cos^2(\theta_{k,t})\Bigr).
\]
则有完全等价的多项式展开
\[
\boxed{\;
\EE\bigl[u^{D_m}\mid \mathcal C_m\bigr]
=\frac{1}{2^m\,\mathsf{Col}_m}\sum_{I\subseteq[m]} E_{m,I}\,(u+1)^{m-|I|}(1-u)^{|I|}.\;}
\]
\end{proposition}

\begin{proof}
由引理 \ref{lem:pom-fold-congruence}，\(\mathcal C_m\) 等价于同余 \(N(\omega)\equiv N(\eta)\pmod M\)。
因此对任意 \(a\in\ZZ/M\ZZ\) 的角色恒等式
\(
\mathbf 1_{\{a=0\}}=\frac1M\sum_{t=0}^{M-1}e^{2\pi i ta/M}
\)
给出
\[
\mathbf 1_{\mathcal C_m}
=\frac1M\sum_{t=0}^{M-1}
e^{\frac{2\pi i t}{M}(N(\omega)-N(\eta))}.
\]
注意到
\(
N(\omega)-N(\eta)=\sum_{k=1}^m(\omega_k-\eta_k)F_{k+1}
\)
且坐标独立，于是
\[
\EE\bigl[u^{D_m}\mathbf 1_{\mathcal C_m}\bigr]
=\frac1M\sum_{t=0}^{M-1}\prod_{k=1}^m
\EE\Bigl[u^{\mathbf 1_{\{\omega_k\neq\eta_k\}}}e^{\frac{2\pi i t}{M}(\omega_k-\eta_k)F_{k+1}}\Bigr].
\]
对固定 \((k,t)\) 枚举 \((\omega_k,\eta_k)\in\{0,1\}^2\) 得内层期望等于
\[
\frac{1}{4}\Bigl(2+u e^{2i\theta_{k,t}}+u e^{-2i\theta_{k,t}}\Bigr)=\frac{1+u\cos(2\theta_{k,t})}{2},
\]
从而得到第一条闭式。
取 \(u=1\) 得碰撞概率的 \(\cos^2\) 乘积表示；再以比值给出条件母函数。
\par
最后用恒等式
\[
\frac{1+u\cos(2\theta)}{2}
=\frac{u+1}{2}\cos^2\theta+\frac{1-u}{2}\sin^2\theta
\]
逐坐标展开乘积，并对 \(t\) 取平均定义 \(E_{m,I}\)，即得多项式形式。证毕。
\end{proof}

\begin{proposition}[碰撞条件下差异场的二阶结构：协方差闭式]\label{prop:fold-collision-conditional-hamming-covariance}
沿用命题 \ref{prop:fold-collision-conditional-hamming-mgf} 的记号。
令 \(\delta_k:=\mathbf 1_{\{\omega_k\neq\eta_k\}}\) 并令
\(
S_k:=(-1)^{\omega_k}(-1)^{\eta_k}\in\{\pm1\}
\)。
则对任意 \(i\neq j\) 有精确协方差公式
\[
\boxed{\;
\mathrm{Cov}(\delta_i,\delta_j\mid \mathcal C_m)
=\frac14\left(\frac{E_{m,\{i,j\}}}{\mathsf{Col}_m}-\frac{E_{m,\{i\}}E_{m,\{j\}}}{\mathsf{Col}_m^{\,2}}\right).\;}
\]
并且条件方差具有闭式分解
\[
\boxed{\;
\mathrm{Var}(D_m\mid\mathcal C_m)
=\frac14\sum_{k=1}^m\left(1-\Bigl(\frac{E_{m,\{k\}}}{\mathsf{Col}_m}\Bigr)^2\right)
\;+\;\frac12\sum_{1\le i<j\le m}\left(\frac{E_{m,\{i,j\}}}{\mathsf{Col}_m}-\frac{E_{m,\{i\}}E_{m,\{j\}}}{\mathsf{Col}_m^{\,2}}\right).\;}
\]
\end{proposition}

\begin{proof}
恒等式 \(\delta_k=\frac{1-S_k}{2}\) 给出
\(
\mathrm{Cov}(\delta_i,\delta_j\mid\mathcal C_m)=\frac14\,\mathrm{Cov}(S_i,S_j\mid\mathcal C_m)
\)。
由命题 \ref{prop:fold-collision-conditional-hamming-mgf} 的 Fourier 表示，内层坐标期望满足
\[
\EE\bigl[S_k\,\mathbf 1_{\mathcal C_m}\bigr]
=\frac1M\sum_{t=0}^{M-1}\sin^2(\theta_{k,t})\prod_{\ell\ne k}\cos^2(\theta_{\ell,t})
=E_{m,\{k\}},
\]
并且类似地
\(
\EE[S_iS_j\,\mathbf 1_{\mathcal C_m}]=E_{m,\{i,j\}}
\)。
除以 \(\PP(\mathcal C_m)=\mathsf{Col}_m\) 得
\(
\EE[S_k\mid\mathcal C_m]=E_{m,\{k\}}/\mathsf{Col}_m
\)、
\(
\EE[S_iS_j\mid\mathcal C_m]=E_{m,\{i,j\}}/\mathsf{Col}_m
\)，
代入协方差定义即得第一式。
\par
对方差分解，使用
\(
D_m=\sum_{k=1}^m\delta_k
\)
与
\(
\mathrm{Var}(\sum_k\delta_k)=\sum_k\mathrm{Var}(\delta_k)+2\sum_{i<j}\mathrm{Cov}(\delta_i,\delta_j)
\)，
并注意 \(\mathrm{Var}(\delta_k\mid\mathcal C_m)=\frac14\bigl(1-\EE[S_k\mid\mathcal C_m]^2\bigr)\)。
代入并整理即得所示闭式。证毕。
\end{proof}

\endinput

